% FAC-EN-U0041
% French source: fac_body.tex lines 1444--1486.
% Authority: PDF physical pp. 39--41; printed pp. 234--236 (no. 43 complete).

\medskip
\phantomsection\label{fac:II:3:no:43}
\textbf{43. Some preliminary properties of irreducible varieties.}
Let $V$ be a closed subvariety of $K^r$, and let $I(V)$ be the ideal of $K[X_1,\ldots,X_r]$ consisting of the polynomials vanishing on $V$. Let
\[
A=K[X_1,\ldots,X_r]/I(V).
\]
There is a canonical homomorphism
\[
\iota:A\longrightarrow\Gamma(V,\mathcal{O}_V),
\]
which is injective by the very definition of $I(V)$.

\phantomsection\label{fac:II:3:no:43:prop:4}
\textsc{Proposition 4.} \textit{If $V$ is irreducible, then $\iota:A\to\Gamma(V,\mathcal{O}_V)$ is bijective.}

(In fact, this holds for \textit{every} closed subvariety of $K^r$, as we shall show in the \hyperref[fac:II:3:no:44]{next number}.)

Let $K(V)$ be the field of fractions of $A$. By \hyperref[fac:II:1:no:36]{no.~36}, $\mathcal{O}_{x,V}$ may be identified with the localization of $A$ at the maximal ideal $\mathfrak{m}_x$ consisting of the polynomials vanishing at $x$, and
\[
\Gamma(V,\mathcal{O}_V)=A=\bigcap_{x\in V}\mathcal{O}_{x,V},
\]
where all the $\mathcal{O}_{x,V}$ are regarded as subrings of $K(V)$. But every maximal ideal of $A$ is one of the $\mathfrak{m}_x$, because $K$ is algebraically closed (Hilbert's Nullstellensatz). It follows immediately (cf.~[8], Chap.~XV, \S5, Theorem~X) that
\[
A=\bigcap_{x\in V}\mathcal{O}_{x,V}=\Gamma(V,\mathcal{O}_V),
\]
as required.

\phantomsection\label{fac:II:3:no:43:prop:5}
\textsc{Proposition 5.} \textit{Let $X$ be an irreducible algebraic variety, let $Q$ be a regular function on $X$, and let $P$ be a regular function on $X_Q$. Then, for every sufficiently large $n$, the rational function $Q^nP$ is regular on all of $X$.}

By the quasi-compactness of $X$, the question is local; by \hyperref[fac:II:3:no:42:thm:1]{Theorem~1}, we may therefore suppose that $X$ is a closed subvariety of $K^r$. The \hyperref[fac:II:3:no:43:prop:4]{preceding proposition} then shows that $Q$ is an element of
\[
A=K[X_1,\ldots,X_r]/I(X).
\]
The hypothesis on $P$ means that, for every point $x\in X_Q$, we can write $P=P_x/Q_x$, with $P_x,Q_x\in A$ and $Q_x(x)\neq0$. Let $\mathfrak{a}$ be the ideal of $A$ generated by the $Q_x$. The zero locus of $\mathfrak{a}$ is contained in the zero locus of $Q$; by Hilbert's Nullstellensatz, it follows that $Q^n\in\mathfrak{a}$ for all sufficiently large $n$. Hence
\[
Q^n=\sum R_xQ_x,
\qquad
Q^nP=\sum R_xP_x,
\qquad R_x\in A,
\]
which shows that $Q^nP$ is regular on $X$.

(One could instead use the fact that $X_Q$ is affine when $X$ is affine and apply \hyperref[fac:II:3:no:43:prop:4]{Proposition~4} to $X_Q$.)

\phantomsection\label{fac:II:3:no:43:prop:6}
\textsc{Proposition 6.} \textit{Let $X$ be an irreducible algebraic variety, let $Q$ be a regular function on $X$, let $\mathcal{F}$ be a coherent algebraic sheaf on $X$, and let $s$ be a section of $\mathcal{F}$ over $X$ whose restriction to $X_Q$ is zero. Then, for every sufficiently large $n$, the section $Q^ns$ is zero on all of $X$.}

Again the question is local, so we may suppose:

\begin{enumerate}
\item[(a)]\phantomsection\label{fac:II:3:no:43:prop:6:item:a} $X$ is a closed subvariety of $K^r$;
\item[(b)]\phantomsection\label{fac:II:3:no:43:prop:6:item:b} $\mathcal{F}$ is isomorphic to the cokernel of a homomorphism $\varphi:\mathcal{O}_X^p\to\mathcal{O}_X^q$;
\item[(c)]\phantomsection\label{fac:II:3:no:43:prop:6:item:c} $s$ is the image of a section $\sigma$ of $\mathcal{O}_X^q$.
\end{enumerate}

(Indeed, all these conditions hold locally.)

Set
\[
A=\Gamma(X,\mathcal{O}_X)=K[X_1,\ldots,X_r]/I(X).
\]
The section $\sigma$ may be identified with a system of $q$ elements of $A$. On the other hand, let
\[
t_1=\varphi(1,0,\ldots,0),\ldots,
t_p=\varphi(0,\ldots,0,1).
\]
\phantomsection\label{fac:II:3:no:43:prop:6:proof:denominator-clearing}
The $t_i$, for $1\leq i\leq p$, are sections of $\mathcal{O}_X^q$ over $X$, and hence may also be identified with systems of $q$ elements of $A$. The hypothesis on $s$ means that, for every $x\in X_Q$, we have $\sigma(x)\in\varphi(\mathcal{O}_{x,X}^p)$; that is, $\sigma$ can be written in the form
\[
\sigma=\sum_{i=1}^{i=p}f_it_i,
\qquad f_i\in\mathcal{O}_{x,X}.
\]
Equivalently, after clearing denominators, there is a $Q_x\in A$, with $Q_x(x)\neq0$, such that
\[
Q_x\sigma=\sum_{i=1}^{i=p}R_it_i,
\qquad R_i\in A.
\]
The \hyperref[fac:II:3:no:43:prop:6:proof:denominator-clearing]{preceding argument} then shows that, for every sufficiently large $n$, $Q^n$ belongs to the ideal generated by the $Q_x$. Consequently $Q^n\sigma(x)\in\varphi(\mathcal{O}_{x,X}^p)$ for every $x\in X$, which says precisely that $Q^ns$ is zero on all of $X$.
